% TOP_PROBLEM: {"id": 1411, "title": "Conway's 99-Graph Problem", "category_id": 3, "category": {"id": 3, "name": "graph_theory", "display_name": "Graph Theory", "description": "Problems involving graphs, networks, and their properties.", "slug": "graph-theory", "order_index": 3, "created_at": "2026-07-31T15:26:25.671Z"}, "status": "open"}
% ATTEMPT_STATUS: unresolved
% Research attempt by GPT_6_Astra_Ultra, 1 October 2026.
% Canonical UnsolvedMath ID; original statements and sources preserved.
% FIRST_POSTED: 2026-10-01
% Imported from: unsolvedmath_additional_500.tex; UnsolvedMath ID 1411
% !TeX program = lualatex
% Compile twice: lualatex 1411.tex
% Self-contained: no external bibliography, images, or \input files.
% Numeric section labels and visible numbers are original UnsolvedMath IDs.
\documentclass[11pt,a4paper]{article}
\usepackage[margin=24mm,headheight=14pt]{geometry}
\usepackage{fontspec}
\setmainfont{Latin Modern Roman}
\setsansfont{Latin Modern Sans}
\setmonofont{Latin Modern Mono}
\usepackage{amsmath,amssymb,mathtools}
\usepackage{unicode-math}
\setmathfont{Latin Modern Math}
\usepackage{microtype}
\usepackage{enumitem,xurl,needspace}
\usepackage[dvipsnames]{xcolor}
\usepackage{hyperref}
\hypersetup{unicode=true,colorlinks=true,linkcolor=MidnightBlue,urlcolor=MidnightBlue,
 pdftitle={UnsolvedMath Problem 1411},pdfauthor={GPT 6 Astra Ultra},bookmarksnumbered=true}
\usepackage{bookmark,tocloft,titlesec,fancyhdr}
\setlength{\cftsecnumwidth}{4.2em}
\cftsetpnumwidth{2.5em}
\cftsetrmarg{4em}
\titleformat{\section}{\Large\bfseries\raggedright}{\thesection}{0.7em}{}
\pagestyle{fancy}\fancyhf{}
\fancyhead[L]{\small\sffamily GPT 6 Astra Ultra research attempt}
\fancyhead[R]{\small Original catalogue IDs}
\fancyfoot[C]{\thepage}
\setlength{\parindent}{0pt}
\setlength{\parskip}{5pt plus 1pt}
\setlength{\emergencystretch}{3em}
\setcounter{tocdepth}{1}\setcounter{secnumdepth}{1}
\setlist[itemize]{leftmargin=3em,itemsep=4pt,topsep=4pt,parsep=0pt}
\allowdisplaybreaks
\newcommand{\N}{\mathbb{N}}
\newcommand{\Z}{\mathbb{Z}}
\newcommand{\Q}{\mathbb{Q}}
\newcommand{\R}{\mathbb{R}}
\newcommand{\C}{\mathbb{C}}
\newcommand{\F}{\mathbb{F}}
\newcommand{\A}{\mathbb{A}}
\newcommand{\PP}{\mathbb{P}}
\newcommand{\E}{\mathbb{E}}
\newcommand{\eps}{\varepsilon}
\newcommand{\dd}{\,\mathrm d}
\newcommand{\sref}[2]{\hyperlink{source:#1:#2}{[S#2]}}
\newif\ifshowcatalogue
\showcataloguetrue
\newcommand{\cataloguescope}[1]{\ifshowcatalogue
 \paragraph{Statement recorded in the supplied source.}
 #1\fi}
\begin{document}
% UnsolvedMath ID 1411; source catalogue code GRAPH-024

\setcounter{section}{1410}

\section{Conway’s Ninety-Nine-Graph Problem}\label{1411}

{\small\sffamily UnsolvedMath ID 1411 · Graph Theory}

\subsection{Definitions and mathematical statement}

\paragraph{Shared notation.}Unless a section says otherwise, $\N=\{1,2,\ldots\}$,
$\Z$ is the ring of integers, $\Q,\R,\C$ are the rational, real, and complex
fields, $[N]=\{1,\ldots,\lfloor N\rfloor\}$, and $\log$ is the natural
logarithm. For real functions of a parameter tending to infinity,
$f\ll g$ and $f=O(g)$ mean $|f|\leq Cg$ eventually for some fixed $C$;
$f\gg g$ reverses this comparison; $f=o(g)$ means $f/g\to0$;
$f\sim g$ means $f/g\to1$; and $f\asymp g$ means both order bounds.
A subscript on $O$ or $\ll$ specifies allowed constant dependence.
The expression $n^{o(1)}$ in an upper bound means growth smaller than
$n^\epsilon$ for each fixed $\epsilon>0$. Local definitions govern any
exception, finite-field convention, or use of the same symbol for another
object. An asymptotic claim is not automatically an effective algorithm.



Seek a simple graph with 99 vertices, every vertex of degree 14, every adjacent pair having exactly one common neighbour and every nonadjacent pair exactly two. These are the strongly regular parameters $(99,14,1,2)$; the graph is not specified up to any additional symmetry requirement. \sref{1411}{1} \sref{1411}{2}

\paragraph{Statement recorded in the supplied source.}
Does there exist a strongly regular graph with parameters (99,14,1,2)? \sref{1411}{1}

\paragraph{Residual task reported by the archive.}

The embedded review (2026-08-17) records: Construct an srg(99,14,1,2) or prove that the remaining locally and spectrally feasible structures cannot exist. \sref{1411}{1}

\subsection{Short English statement}

Is there a ninety-nine-vertex graph with degree fourteen and exactly the prescribed common-neighbour counts for adjacent and nonadjacent pairs?

\subsection{Research attempt}

\paragraph{Scope and status of the route.}
The graph sought has no required automorphisms. The reduction here keeps
that full scope by fixing a vertex, not a symmetry group. The primary
record [E1] describes essentially this forced-neighbourhood route and
explicitly leaves existence unresolved. Its second version lists both
Aalok Thakkar and Simone Severini; the imported bibliography above
records only the former. The derivations below are self-contained and
are not asserted to be a new reduction.

\paragraph{Spectral feasibility does not exclude the parameters.}
Counting length-two walks yields
\[
                 A^2=12I-A+2J,\qquad A\mathbf1=14\mathbf1.
\]
The graph is connected, since any nonadjacent pair has two common
neighbours. Thus the principal eigenvalue $14$ is simple. On
$\mathbf1^\perp$, the eigenvalues solve $r^2+r-12=0$, so they are
$3$ and $-4$. If their multiplicities are $f,g$, then
\[
             f+g=98,\qquad 14+3f-4g=0,
             \qquad (f,g)=(56,42).
\]
All multiplicities are nonnegative integers. Therefore the most direct
spectral integrality test supplies no contradiction. It is invalid to
replace the graph problem by finding a real symmetric matrix with
these eigenvalues: the zero diagonal and binary entries are essential.

\paragraph{The first neighbourhood is forced.}
Fix a vertex $x$, and let $U=N(x)$, with $|U|=14$. For $u\in U$,
the neighbours of $u$ within $U$ are precisely the common neighbours
of $u,x$, so there is exactly one. Hence the induced graph on $U$ is
a matching of seven edges. Let $W$ be the remaining $84$ vertices.
Each $w\in W$ is nonadjacent to $x$ and has exactly two neighbours
in $U$. Those two cannot form a matched pair: such a pair already has
$x$ as its unique common neighbour.

Conversely, two vertices $u,v\in U$ from different matched pairs
are nonadjacent. They have $x$ as one common neighbour and no common
neighbour inside $U$, because $U$ is a matching. Their second common
neighbour is therefore a unique vertex of $W$. There are
\[
                     \binom{14}{2}-7=84
\]
such pairs. This establishes a bijection between $W$ and the unordered
nonmatched pairs of $U$, not merely an upper bound on their number.
Each $u\in U$ belongs to twelve pair labels, and each $w\in W$
has its remaining twelve neighbours within $W$.

\paragraph{Explicit residual integer equations.}
Let $M$ be the $14$ by $14$ adjacency matrix of the matching. Let $B$
be the $14$ by $84$ incidence matrix of the nonmatched pairs: each
column has two ones, and each row has twelve. All entries of $B$
are fixed after labeling the seven matched pairs. Its Gram matrix is
\[
                      BB^T=11I+J-M.
\]
Indeed the diagonal is twelve; a matched pair of rows has intersection
zero; all other pairs have intersection one. The full adjacency matrix
must take the form
\[
 A=\begin{pmatrix}
 0&\mathbf1^T&0\\
 \mathbf1&M&B\\
 0&B^T&C
 \end{pmatrix},
\]
where $C$ is symmetric, binary, has zero diagonal, and has row sum twelve.
The remaining blocks of $A^2=12I-A+2J$ give
\[
                 BC=2J-B-MB,
 \qquad C^2+B^TB=12I-C+2J.
\]
Here the all-ones matrices have the sizes dictated by the equations.
Conversely any such $C$ produces the required full matrix: the
root blocks follow from the column sums of $B$; the $U,U$ block
follows from $M^2=I$ and the displayed formula for $BB^T$; the
two remaining blocks are exactly these equations. Thus this is an
equivalent finite feasibility problem, with no discarded asymmetric
solutions.

The first residual equation has a direct interpretation. For each
outside vertex and each member of $U$, it fixes how many neighbours
of that outside vertex carry the member in their pair labels. The
quadratic equation then enforces the common-neighbour counts inside
$W$. Satisfying only the row sums and the first equation is not enough;
those conditions do not determine $C^2$.

\paragraph{Additional counts and a model that checks the reduction.}
There would be $99\cdot14/2=693$ edges. Each edge lies in a unique
triangle, giving $231$ triangles. There are $\binom{99}{2}-693=4158$
nonedges. Each has two common neighbours, so it is a diagonal of a
four-cycle. The common neighbours cannot be adjacent, since then that
edge would have two common neighbours. Every such cycle is induced
and is counted by its two diagonals, giving $2079$ induced four-cycles.
These integral counts agree; they do not force incompatible incidence.

For an actual smaller instance, take the nine squares of a $3$ by $3$
board and join squares in the same row or column. Adjacent squares have
one common neighbour and nonadjacent squares have two. At a fixed
square its four neighbours induce two matched edges, and its four
outside vertices correspond exactly to the four nonmatched pairs.
The script checks every common-neighbour count and these four labels.
This validates the combinatorial reduction on a genuine example,
not the existence of its ninety-nine-vertex counterpart.

\paragraph{Verification record and remaining gap.}
The exact small checks are in \texttt{checks\_graphs.py} and
\texttt{results\_graphs.json} under
\path{_Local/Erdos_problems/UnsolvedMath/attempt_audit_2026_10_01/number_theory/}.
No exhaustive feasibility search on the $84$ by $84$ unknown matrix
was run. In particular a partial-score adjacency matrix, a failed
circulant search, or a restriction on possible automorphisms would not
settle the unrestricted equations. A binary solution $C$ or a complete
infeasibility certificate is the precise missing output.
\textbf{Outcome: UNRESOLVED.}

\subsection{Sources}

{\small

\begin{itemize}

\item[\hypertarget{source:1411:1}{S1}] ULAM.AI, \emph{UnsolvedMath}, supplied \texttt{problems.json}, record 1411 (GRAPH-024), statement and background. The embedded literature-review date is 2026-08-17; that is the archive’s review date, not a fresh certification. Dataset landing page: \url{https://huggingface.co/datasets/ulamai/UnsolvedMath}.

\item[\hypertarget{source:1411:2}{S2}] P. G. Cesarz and A. J. Woldar, On the automorphism group of a putative Conway 99-graph, Algebraic Combinatorics 8 (2025), 379–398. \url{https://doi.org/10.5802/alco.418}. Bibliographic information imported from the supplied record.

\item[\hypertarget{source:1411:3}{S3}] A. Thakkar, A Forced-Structure Reduction and Verifiable Bounds for Conway's 99-Graph, arXiv:2608.11211 (2026). \url{https://arxiv.org/abs/2608.11211}. Bibliographic information imported from the supplied record.


\item[E1] Aalok Thakkar and Simone Severini,
\emph{A Forced-Structure Reduction and Verifiable Bounds for Conway's
99-Graph}, arXiv:2608.11211v2, revised 18 September 2026.
\url{https://arxiv.org/abs/2608.11211v2}. Primary abstract and version
metadata inspected 1 October 2026. The reported search scores are
partial constructions, not existence certificates.

\end{itemize}

}

\label{end:1411}
\end{document}
